% !TeX root = ../Lua-document-Template.tex %main.tex
% !TeX spellcheck = uk_UA

\setcounter{secnumdepth}{-1} % Тимчасово вимикаємо нумерацію розділів
\chapter{ВСТУП}             % Створюємо розділ (він автоматично піде в зміст і закладки PDF)
\setcounter{secnumdepth}{2}  % Повертаємо стандартну нумерацію назад для наступних глав

\textbf{Актуальність теми.} Шаруваті матеріали широко використовуються в машинобудуванні, авіації, будівництві та електроніці, адже поєднання шарів з різними властивостями дає змогу створювати міцні й водночас легкі конструкції. Проте на межі між шарами з різною жорсткістю виникає концентрація напружень, яка може призвести до розшарування та руйнування. Тому важливо вміти визначати, як розподіляються переміщення і напруження в таких тілах. Задача повздовжнього зсуву є однією з найпростіших моделей такого типу: вона зводиться до рівняння Лапласа в кожному шарі з умовами контакту на межі між ними, що дозволяє отримати аналітичний розв’язок. В останні роки дослідження задач антиплоскої деформації шаруватих тіл здебільшого спрямовані на вивчення тріщин на межі поділу шарів, функціонально-градієнтних матеріалів та застосування чисельних методів, зокрема методу скінченних елементів. Водночас чисельні підходи не дають явної залежності розв’язку від параметрів задачі й потребують перевірки на точних розв’язках, а відомі аналітичні результати переважно стосуються однорідних або необмежених областей. Тому побудова аналітичного розв’язку для обмеженої двошарової області з мішаними граничними умовами, корисного як для інженерних розрахунків, так і для перевірки чисельних методів, є актуальною, особливо з огляду на постійне зростання використання композиційних матеріалів у сучасній техніці.

\textbf{Мета роботи} — побудувати та дослідити розв’язок задачі повздовжнього зсуву для двошарової прямокутної області, на верхній межі якої діє дотичне навантаження $p(x)$, із застосуванням методу інтегральних перетворень.

Для досягнення мети необхідно:
\begin{itemize}
	\item сформулювати математичну постановку задачі з граничними умовами та умовами контакту між шарами;
	\item застосувати скінченне інтегральне перетворення Фур’є за змінною $x$ і звести задачу до системи звичайних диференціальних рівнянь;
	\item розв’язати отриману задачу в просторі зображень і повернутися до оригіналів за допомогою формули оберненого перетворення;
	\item знайти переміщення та дотичні напруження в кожному шарі й проаналізувати вплив параметрів задачі на отримані результати.
\end{itemize}

\textbf{Об’єкт дослідження} — повздовжня зсувна деформація двошарового пружного тіла прямокутної форми.

\textbf{Предмет дослідження} — математична модель задачі повздовжнього зсуву для двошарової прямокутної області, її розв’язання методом інтегральних перетворень, а також розподіл переміщень і напружень у шарах.
